% Options for packages loaded elsewhere
\PassOptionsToPackage{unicode}{hyperref}
\PassOptionsToPackage{hyphens}{url}
\documentclass[
  ignorenonframetext,
]{beamer}
\newif\ifbibliography
\usepackage{pgfpages}
\setbeamertemplate{caption}[numbered]
\setbeamertemplate{caption label separator}{: }
\setbeamercolor{caption name}{fg=normal text.fg}
\beamertemplatenavigationsymbolsempty
% remove section numbering
\setbeamertemplate{part page}{
  \centering
  \begin{beamercolorbox}[sep=16pt,center]{part title}
    \usebeamerfont{part title}\insertpart\par
  \end{beamercolorbox}
}
\setbeamertemplate{section page}{
  \centering
  \begin{beamercolorbox}[sep=12pt,center]{section title}
    \usebeamerfont{section title}\insertsection\par
  \end{beamercolorbox}
}
\setbeamertemplate{subsection page}{
  \centering
  \begin{beamercolorbox}[sep=8pt,center]{subsection title}
    \usebeamerfont{subsection title}\insertsubsection\par
  \end{beamercolorbox}
}
% Prevent slide breaks in the middle of a paragraph
\widowpenalties 1 10000
\raggedbottom
\AtBeginPart{
  \frame{\partpage}
}
\AtBeginSection{
  \ifbibliography
  \else
    \frame{\sectionpage}
  \fi
}
\AtBeginSubsection{
  \frame{\subsectionpage}
}
\usepackage{iftex}
\ifPDFTeX
  \usepackage[T1]{fontenc}
  \usepackage[utf8]{inputenc}
  \usepackage{textcomp} % provide euro and other symbols
\else % if luatex or xetex
  \usepackage{unicode-math} % this also loads fontspec
  \defaultfontfeatures{Scale=MatchLowercase}
  \defaultfontfeatures[\rmfamily]{Ligatures=TeX,Scale=1}
\fi
\usepackage{lmodern}
\ifPDFTeX\else
  % xetex/luatex font selection
\fi
% Use upquote if available, for straight quotes in verbatim environments
\IfFileExists{upquote.sty}{\usepackage{upquote}}{}
\IfFileExists{microtype.sty}{% use microtype if available
  \usepackage[]{microtype}
  \UseMicrotypeSet[protrusion]{basicmath} % disable protrusion for tt fonts
}{}
\makeatletter
\@ifundefined{KOMAClassName}{% if non-KOMA class
  \IfFileExists{parskip.sty}{%
    \usepackage{parskip}
  }{% else
    \setlength{\parindent}{0pt}
    \setlength{\parskip}{6pt plus 2pt minus 1pt}}
}{% if KOMA class
  \KOMAoptions{parskip=half}}
\makeatother
\setlength{\emergencystretch}{3em} % prevent overfull lines
\providecommand{\tightlist}{%
  \setlength{\itemsep}{0pt}\setlength{\parskip}{0pt}}
% Slide layout defaults for warning-clean scientific decks.
\usepackage{etoolbox}
\IfFileExists{xurl.sty}{\usepackage{xurl}}{}
\IfFileExists{seqsplit.sty}{\usepackage{seqsplit}}{\newcommand{\seqsplit}[1]{#1}}
\protected\def\breakseq#1{\seqsplit{#1}}
\protected\def\breaktt#1{\begingroup\ttfamily\seqsplit{#1}\endgroup}
\setlength{\emergencystretch}{6em}
\tolerance=5000
\hbadness=10000
\hfuzz=1pt
\setlength{\tabcolsep}{2pt}
\AtBeginEnvironment{longtable}{\tiny\renewcommand{\arraystretch}{0.86}\setlength{\tabcolsep}{1pt}}
\AtBeginEnvironment{tabular}{\tiny\renewcommand{\arraystretch}{0.86}\setlength{\tabcolsep}{1pt}}

% Natbib and cross-reference fallbacks — slides don't load natbib
% or cleveref, but combined-PDF manuscript prose may emit these
% commands. The fallback renders citations as a bracketed key list
% and cross-references as detokenized labels so slides don't fail on
% undefined control sequences or raw underscores. \providecommand is
% a no-op if packages load later.
\providecommand{\citep}[1]{[#1]}
\providecommand{\citet}[1]{#1}
\providecommand{\citealp}[1]{#1}
\providecommand{\citeauthor}[1]{#1}
\providecommand{\citeyear}[1]{#1}
\providecommand{\cref}[1]{\texttt{\detokenize{#1}}}
\providecommand{\Cref}[1]{\texttt{\detokenize{#1}}}
\usepackage{bookmark}
\IfFileExists{xurl.sty}{\usepackage{xurl}}{} % add URL line breaks if available
\urlstyle{same}
% fallback for those not using the hyperref driver hyperxmp:
\makeatletter
\@ifundefined{xmpquote}{\newcommand{\xmpquote}[1]{#1}}{}
\makeatother
\hypersetup{
  hidelinks,
  pdfcreator={LaTeX via pandoc}}

\author{\texorpdfstring{}{}}
\date{}

\begin{document}

\section{Engine Architecture}\label{engine-architecture}

\begin{frame}[fragile,allowframebreaks]{Engine Architecture}
The engine is nine small modules, each with one responsibility.

\texttt{geometry} is pure arithmetic over points: the page trim, its
margins, and a \texttt{ColumnGrid} that partitions a region into equal
columns with gutters and reports the centre of each gutter (where a
hairline rule is drawn). It imports no rendering library, which makes
the column mathematics trivially testable.

\texttt{content} defines the typed object graph --- \texttt{Edition},
\texttt{Page}, \texttt{Story}, \texttt{Box}, \texttt{Figure},
\texttt{Block} --- and the strict YAML loaders that build it. The
loaders are forgiving about \emph{shape} (a body element may be a bare
string or a tagged mapping) but strict about \emph{structure}: an
unknown page template or box kind raises an error naming the offending
value and the allowed set.

\texttt{config} is the strict render configuration (trim size, default
column count, rail width, gutter), validated the same way.

\texttt{typography} owns every font name: it registers the Didot
display, Georgia text and Helvetica label faces (with a base-14
fall-back when the system lacks them) and builds the paragraph
stylesheet that the rest of the engine draws from through named
\texttt{Fonts} roles.

\texttt{figures} is the only module that \emph{generates} imagery: the
halftone harbour, lighthouse and redwood engravings (Pillow), the
weather, tide and crab-landings charts (Matplotlib), and the colour
display-ad graphics. It writes deterministic PNGs that
\texttt{components} later places.

\texttt{components} builds the flowables that live \emph{inside} a
column: a story's headline, byline and body; a raised drop cap; a pull
quote; a ruled box; a data table; a figure with its cutline.
\texttt{furniture} draws the elements that are \emph{placed} rather than
flowed: the nameplate and its ears, the section banner, the folio, the
hairline column rules, and the spanning lead headline.

\texttt{layout} is the bridge. For each page it draws the furniture
(which returns the y-coordinate where the column grid begins),
constructs the column frames --- including an optional narrow rail and
the room reserved for a spanning headline --- and pours the flowables
into the frames with ReportLab's \breaktt{Frame.addFromList}, which
splits paragraphs across columns automatically. Content that does not
fit is returned and reported as an \emph{over-set} page rather than
silently dropped.

\texttt{engine} ties it together: register fonts, build the stylesheet,
open a canvas at the configured trim, render each page, and write the
PDF plus a machine-readable render report.

This \textbf{drawn-furniture / flowed-body} split is the key
architectural choice. A pure document-template approach cannot place a
headline that spans several columns above body text that flows beneath
it; drawing the furniture first, and beginning the column frames below
it, makes spanning headlines, standing boxes and automatic text flow
coexist on the same page.
\end{frame}

\begin{frame}[fragile,allowframebreaks]{Methods}
\protect\phantomsection\label{methods}
Producing an edition is a deterministic, four-step method, run by the
pipeline and covered end-to-end by the test suite:

\begin{enumerate}
\tightlist
\item
  \textbf{Load} --- \texttt{content} parses \texttt{edition.yaml} and
  every \breaktt{content/pages/*.yaml} into the typed \texttt{Edition}
  graph, rejecting unknown page templates or box kinds by naming the
  offending value and the allowed set.
\item
  \textbf{Measure} --- \texttt{geometry} derives the trim, margins and
  \texttt{ColumnGrid} for each page from the strict \texttt{config},
  independently of any rendering library.
\item
  \textbf{Compose} --- \texttt{components} and \texttt{furniture} build
  the flowed and drawn elements; \texttt{layout} pours them into the
  column frames and reports any content that does not fit as an over-set
  page rather than dropping it silently.
\item
  \textbf{Emit} --- \texttt{engine} registers fonts, renders each page
  to the canvas, and writes the PDF alongside a machine-readable render
  report.
\end{enumerate}

Every step is pure-data-in, artifact-out and seeded deterministically,
so the same edition manifest always yields a byte-identical paper ---
the reproducibility contract verified in \texttt{tests/} and described
in \breaktt{04\_reproducibility.md}.
\end{frame}

\end{document}
